\section{Conclusion}\label{sec:conclusion}
Neural Bellman Operators organize policy evaluation, feasible improvement,
and economic assessment around a specified continuation. The contribution
is an economic computation with an explicit error and work account, not a
new principle of optimality or a presumption that a neural approximation
must dominate a conventional one.

The centered decision-loss bound identifies the errors that can change
current choices. The continuation-transport theorem extends reuse to a
verified neighborhood of future economic primitives and gives a condition
for refreshing the evaluated object. The scalar true-objective certificate
then supports actual first-success stopping with every failed check charged.
These results connect approximation, economic stability, verification, and
amortization without conflating their distinct assumptions.

The capital evidence contains both successful NBO decisions and useful
reversals. The original high-dimensional menus establish protected payoff
advantages in declared cells, while other cells favor simulation and simpler
actors. The new query-volume study certifies NBO across every declared
stream and retains the competitive conventional surrogates, unresolved actor
outcomes, full verification bill, and own-action diagnostics. A repetition
of the same frozen study is a replication of its computation, not a larger
independent experiment. The finite scalar certificate does not substitute
for the separate continuous-time and full-control accounts.

The future-change experiment makes refresh an executed decision rather than
an informal qualification. NBO certifies the full scalar interval after every
registered refresh, and its action certificates establish economically signed
responses to changes in future production and valuation. The complete work
record shows the additional price of failed reuse checks and leaves the
conventional comparators on the frontier. The policy-composition theorem
connects this local verification architecture to a full dynamic target when
own-policy continuation and uniform coverage are supplied. It does not infer
those missing global premises from a finite mean-loss table.

Recursive utility, endogenous preferences, temporal selves, and dynamic
games remain part of the paper. Their utility domains, information sets,
transversality conditions, and deviation criteria determine the continuations
that can be evaluated and reused. The complete application models and proofs
are retained. Taken together, the results specify when a learned scalar
continuation supports an economically accurate decision, when transport or
refresh is required, and how to compare its complete computational cost with
other legitimate numerical procedures.
